\section{Operational adequacy for call-by-name PCF}
\label{sec:p10-operational-adequacy-domain-models}

A continuous interpretation is not automatically an adequate programming
model. We now specify the language, evaluation and observations before
proving their agreement with the domain interpretation. The language is
Gunter's author-presentation PCF, including its predecessor-at-zero
convention \cite[Sections~4.1--4.4, pp.~29--35]{Gunter1994TypesAuthorOA}.
The adequacy argument supplies the proof omitted at that source's p.~35;
it uses a typed logical relation and the least-fixed-point construction,
not the free-powerdomain property. Plotkin's original computability
argument is the primary proof-method reference
\cite[Theorem~3.1; Lemmas~3.2--3.3, pp.~230--232]{Plotkin1977LCFAuthorOA}.
His original constants and predecessor convention are not silently
identified with this Gunter language. Abramsky--Jung expressly separate
the mathematical domain theory from detailed applications
\cite[pp.~7,148--149]{AbramskyJungCorrectedExpandedDomainTheory}.

\subsection{The language, strategy and observations}

\begin{hbdefinition}{def:p10-oa-language}{Typed PCF with the source's recursion binder}
Types and terms are
\[
 t::=\mathsf{num}\mid\mathsf{bool}\mid t\to t,
\]
\[
 \begin{split}
 M::={}&0\mid\mathsf{true}\mid\mathsf{false}
 \mid\operatorname{succ}(M)\mid\operatorname{pred}(M)
 \mid\operatorname{zero?}(M)\\
 &\mid\mathsf{if}\ M\ \mathsf{then}\ M\ \mathsf{else}\ M
 \mid x\mid\lambda x:t.M\mid M\,M\mid\mu x:t.M.
 \end{split}
\]
A finite context \(H\) assigns types to distinct variables.
The source judgment \(H\vdash M:t\) has variable, application and
abstraction rules, the displayed ground constructors, conditional with
same-typed branches, and recursion: the body and bound variable of
\(\mu x:t.M\) both have type \(t\). This \(\mu\) is a term binder,
not a probability mean or a semantic monad multiplication.
Both binders use capture-avoiding substitution; write \([N/x]M\).
Let \(\bar n=\operatorname{succ}^n(0)\).
The closed well-typed values are numerals, the two booleans and typed
abstractions. Function values need not have terminating bodies.
\cite[pp.~17--18,29--31]{Gunter1994TypesAuthorOA}.
\end{hbdefinition}

\begin{hbdefinition}{def:p10-oa-evaluation}{Weak call-by-name evaluation and ground results}
Write \(M\to N\) for the small-step relation registered below.
The argument of a function is substituted without being evaluated.
No reduction is performed under \(\lambda\), in an argument position,
or in an unselected conditional branch. Evaluation contexts are
\[
 E::=[]\mid\operatorname{succ}(E)\mid\operatorname{pred}(E)
 \mid\operatorname{zero?}(E)\mid E\,L
 \mid\mathsf{if}\ E\ \mathsf{then}\ M\ \mathsf{else}\ N.
\]
They propagate the registered primitive contractions. In particular
\(\operatorname{pred}(0)\) contracts to \(0\), and the recursion
binder unfolds once by substituting itself into its body.
Write \(M\Downarrow V\) for the least relation generated by the
registered natural-evaluation rules, and \(\to^*\) for the finite
reflexive transitive closure. A program is a closed term of ground type.
Its observation is its resulting numeral or boolean, when one exists.
Write \(M\Uparrow\) when no such finite evaluation exists.
The symbols \(\to\) and \(\Downarrow\) here describe operational
steps and evaluation, not denotational order or logical implication
\cite[Tables~7--8; Theorem~4.2.2, pp.~31--33]{Gunter1994TypesAuthorOA}.
\end{hbdefinition}

\calculussheet{sheet:p10-operational-model}{Gunter PCF: typing and call-by-name evaluation}{Simple num/bool/arrow types; source mu binder; weak call-by-name steps and natural evaluation; pred(0)=0.}

\hbrule{rule:p10-oa-ty-var}{Var}{}{H\vdash x:t}{\(x:t\in H\).}
\hbrule{rule:p10-oa-ty-abs}{Abs}{H,x:s\vdash M:t}{H\vdash\lambda x:s.M:s\to t}{\(x\text{ fresh for }H\).}
\hbrule{rule:p10-oa-ty-app}{App}{H\vdash M:s\to t\quad H\vdash N:s}{H\vdash M\,N:t}{\none}
\hbrule{rule:p10-oa-ty-zero}{Zero}{}{H\vdash0:\mathsf{num}}{\none}
\hbrule{rule:p10-oa-ty-true}{True}{}{H\vdash\mathsf{true}:\mathsf{bool}}{\none}
\hbrule{rule:p10-oa-ty-false}{False}{}{H\vdash\mathsf{false}:\mathsf{bool}}{\none}
\hbrule{rule:p10-oa-ty-succ}{Succ}{H\vdash M:\mathsf{num}}{H\vdash\operatorname{succ}(M):\mathsf{num}}{\none}
\hbrule{rule:p10-oa-ty-pred}{Pred}{H\vdash M:\mathsf{num}}{H\vdash\operatorname{pred}(M):\mathsf{num}}{\none}
\hbrule{rule:p10-oa-ty-test}{IsZero}{H\vdash M:\mathsf{num}}{H\vdash\operatorname{zero?}(M):\mathsf{bool}}{\none}
\hbrule{rule:p10-oa-ty-if}{Cond}{H\vdash L:\mathsf{bool}\quad H\vdash M:t\quad H\vdash N:t}{H\vdash\mathsf{if}\ L\ \mathsf{then}\ M\ \mathsf{else}\ N:t}{\none}
\hbrule{rule:p10-oa-ty-rec}{Rec}{H,x:t\vdash M:t}{H\vdash\mu x:t.M:t}{\(x\text{ fresh for }H\).}

\hbrule{rule:p10-oa-step-beta}{Beta}{}{(\lambda x:s.M)N\to[N/x]M}{Capture-avoiding substitution; the argument is not evaluated.}
\hbrule{rule:p10-oa-step-pred-zero}{PredZero}{}{\operatorname{pred}(0)\to0}{\none}
\hbrule{rule:p10-oa-step-pred-succ}{PredSucc}{}{\operatorname{pred}(\operatorname{succ}(\bar n))\to\bar n}{\(n\geq0\).}
\hbrule{rule:p10-oa-step-test-zero}{TestZero}{}{\operatorname{zero?}(0)\to\mathsf{true}}{\none}
\hbrule{rule:p10-oa-step-test-succ}{TestSucc}{}{\operatorname{zero?}(\operatorname{succ}(\bar n))\to\mathsf{false}}{\(n\geq0\).}
\hbrule{rule:p10-oa-step-if-true}{IfTrue}{}{\mathsf{if}\ \mathsf{true}\ \mathsf{then}\ M\ \mathsf{else}\ N\to M}{\none}
\hbrule{rule:p10-oa-step-if-false}{IfFalse}{}{\mathsf{if}\ \mathsf{false}\ \mathsf{then}\ M\ \mathsf{else}\ N\to N}{\none}
\hbrule{rule:p10-oa-step-rec}{Unfold}{}{\mu x:t.M\to[\mu x:t.M/x]M}{Capture-avoiding substitution.}
\hbrule{rule:p10-oa-step-context}{EvalContext}{M\to N}{E[M]\to E[N]}{\(E\) is a typed evaluation context of Definition~\ref{def:p10-oa-evaluation}.}

\hbrule{rule:p10-oa-eval-zero}{ZeroValue}{}{0\Downarrow0}{\none}
\hbrule{rule:p10-oa-eval-true}{TrueValue}{}{\mathsf{true}\Downarrow\mathsf{true}}{\none}
\hbrule{rule:p10-oa-eval-false}{FalseValue}{}{\mathsf{false}\Downarrow\mathsf{false}}{\none}
\hbrule{rule:p10-oa-eval-abs}{AbsValue}{}{\lambda x:s.M\Downarrow\lambda x:s.M}{The abstraction is closed and well typed.}
\hbrule{rule:p10-oa-eval-succ}{SuccEval}{M\Downarrow\bar n}{\operatorname{succ}(M)\Downarrow\operatorname{succ}(\bar n)}{\(n\geq0\).}
\hbrule{rule:p10-oa-eval-pred-zero}{PredZeroEval}{M\Downarrow0}{\operatorname{pred}(M)\Downarrow0}{\none}
\hbrule{rule:p10-oa-eval-pred-succ}{PredSuccEval}{M\Downarrow\operatorname{succ}(\bar n)}{\operatorname{pred}(M)\Downarrow\bar n}{\(n\geq0\).}
\hbrule{rule:p10-oa-eval-test-zero}{TestZeroEval}{M\Downarrow0}{\operatorname{zero?}(M)\Downarrow\mathsf{true}}{\none}
\hbrule{rule:p10-oa-eval-test-succ}{TestSuccEval}{M\Downarrow\operatorname{succ}(\bar n)}{\operatorname{zero?}(M)\Downarrow\mathsf{false}}{\(n\geq0\).}
\hbrule{rule:p10-oa-eval-app}{NameApp}{M\Downarrow\lambda x:s.L\quad[N/x]L\Downarrow V}{M\,N\Downarrow V}{No evaluation premise for \(N\).}
\hbrule{rule:p10-oa-eval-if-true}{TrueBranch}{L\Downarrow\mathsf{true}\quad M\Downarrow V}{\mathsf{if}\ L\ \mathsf{then}\ M\ \mathsf{else}\ N\Downarrow V}{No premise for the unselected branch.}
\hbrule{rule:p10-oa-eval-if-false}{FalseBranch}{L\Downarrow\mathsf{false}\quad N\Downarrow V}{\mathsf{if}\ L\ \mathsf{then}\ M\ \mathsf{else}\ N\Downarrow V}{No premise for the unselected branch.}
\hbrule{rule:p10-oa-eval-rec}{RecEval}{[\mu x:t.M/x]M\Downarrow V}{\mu x:t.M\Downarrow V}{The recursive term is substituted, not evaluated as a premise.}


\begin{hbtheorem}{thm:p10-oa-operational}{proved=annex:p10.oa.operational}{Substitution, preservation, progress and the two evaluations}
Typed capture-avoiding substitution preserves typing. Small steps preserve
typing; a closed well-typed nonvalue has exactly one next step, and a
value has none. For closed well-typed terms,
\[
 M\Downarrow V\quad\Longleftrightarrow\quad M\to^*V.
\]
The result is unique up to consistent bound-variable renaming.
At ground type, absence of a finite result is equivalent to an infinite
sequence of the deterministic evaluation steps; it is not a stuck
type error. Moreover \(M\to N\) implies
\(M\Downarrow V\) iff \(N\Downarrow V\) for every value \(V\).
\end{hbtheorem}

\subsection{The actual domain interpretation}

\begin{hbdefinition}{def:p10-oa-interpretation}{Flat observations and full continuous function spaces}
Put \(D_t=C\llbracket t\rrbracket\), using Gunter's interpretation
\(C\). Specifically
\[
 D_{\mathsf{num}}=\mathbb N_{\bot},\quad
 D_{\mathsf{bool}}=\{\bot,\mathsf{true},\mathsf{false}\},\quad
 D_{s\to t}=[D_s\longrightarrow D_t].
\]
Both ground orders are flat, and function spaces have pointwise order.
An \(H\)-environment \(\rho\) supplies \(\rho(x)\in D_{H(x)}\).
Abbreviate \(C\llbracket H\triangleright M:t\rrbracket\rho\) by
\(\llbracket M\rrbracket_\rho\); this drops redundant typing data,
not a source version. Variables select their environment components,
application is evaluation, and abstraction is
\(d\mapsto\llbracket M\rrbracket_{\rho[x\mapsto d]}\).
Ground constants have their corresponding nonbottom values.
Successor, predecessor and zero-test are strict on bottom and have
their ordinary displayed numeral/boolean results, with
\(\operatorname{pred}(0)=0\).
The conditional is strict in its guard, selects its true/false branch,
and returns \(\bot_t\) for bottom guard; it is not strict in both
branches. Finally
\[
 \llbracket\mu x:t.M\rrbracket_\rho
 =\operatorname{fix}\bigl(d\mapsto
                    \llbracket M\rrbracket_{\rho[x\mapsto d]}\bigr).
\]
The function space is not lifted again: this is the call-by-name
interpretation, not a value/computation model for call-by-value
\cite[pp.~34--35]{Gunter1994TypesAuthorOA}.
\end{hbdefinition}

\begin{hbtheorem}{thm:p10-oa-semantic-soundness}{proved=annex:p10.oa.semantic-soundness}{Well-defined continuous meanings and evaluation soundness}
Every \(D_t\) is a pointed bounded-complete algebraic domain.
For each typed term its interpretation is jointly Scott-continuous in
its finite environment. Interpretation respects bound renaming and
typed simultaneous substitution:
\[
 \llbracket M[\sigma]\rrbracket_\rho
   =\llbracket M\rrbracket_{x\mapsto\llbracket\sigma(x)\rrbracket_\rho}.
\]
Every small step preserves denotation. Consequently
\(M\Downarrow V\) implies \(\llbracket M\rrbracket=\llbracket V\rrbracket\).
This proves one direction only; a separate logical-relation argument
is needed for the converse.
\end{hbtheorem}

\subsection{The logical relation and fixed-point induction}

\begin{hbdefinition}{def:p10-oa-logical-relation}{Semantic information related to an operational term}
Relate \(d\in D_t\) to a closed term \(M:t\) by \(d\mathrel{\mathcal R_t}M\).
At ground type, bottom is related to every term; a nonbottom ground
value is related to \(M\) precisely when \(M\) evaluates to its
corresponding numeral or boolean. At arrow type put
\[
 f\mathrel{\mathcal R_{s\to t}}M\quad\Longleftrightarrow\quad
 (\forall N:s\text{ closed})(\forall a\in D_s)\
 a\mathrel{\mathcal R_s}N\ \Longrightarrow\
 f(a)\mathrel{\mathcal R_t}(M\,N).
\]
This relation is between a semantic element and syntax, not equality
of two denotations. Its arrow clause intentionally does not require
\(M\) itself to evaluate to an abstraction.
\end{hbdefinition}

\begin{hbtheorem}{thm:p10-oa-relation-closure}{proved=annex:p10.oa.relation-closure}{Bottom, directed admissibility and operational conversion}
For every closed \(M:t\), the set
\(\{d:d\mathrel{\mathcal R_t}M\}\) contains \(\bot_t\) and is
closed under directed suprema. If \(M\to^*N\), then for every
\(d\in D_t\),
\(d\mathrel{\mathcal R_t}M\) iff \(d\mathrel{\mathcal R_t}N\).
These statements hold at all simple types, including higher arrows.
\end{hbtheorem}

\begin{hbtheorem}{thm:p10-oa-fundamental}{proved=annex:p10.oa.fundamental}{Fundamental lemma, including arbitrary typed recursion}
Let \(H\vdash M:t\). Suppose \(\sigma\) replaces each free
\(x:s\) in \(H\) by a closed term \(\sigma(x):s\), and
\(\rho(x)\mathrel{\mathcal R_s}\sigma(x)\) for every such variable.
Then
\[
                 \llbracket M\rrbracket_\rho
                       \mathrel{\mathcal R_t}M[\sigma].
\]
For a recursion binder, every semantic iterate from bottom is related
to the same full recursive term, by induction on the iterate; directed
admissibility then relates their supremum. Neither that induction nor
its substitution into the body is replaced by a mere fixed-point name.
\end{hbtheorem}

\begin{hbtheorem}{thm:p10-oa-adequacy}{proved=annex:p10.oa.adequacy}{Both directions of ground-observation adequacy}
For every closed program \(M:g\), with
\(g\in\{\mathsf{num},\mathsf{bool}\}\), and every value \(V:g\),
\[
 M\Downarrow V\quad\Longleftrightarrow\quad
                  \llbracket M\rrbracket=\llbracket V\rrbracket.
\]
Consequently \(M\Uparrow\) iff \(\llbracket M\rrbracket=\bot_g\).
This conclusion uses the flat ground carriers and their distinct
nonbottom observations. It is not asserted as termination iff nonbottom
at every higher type.
\end{hbtheorem}

\begin{hbdefinition}{def:p10-oa-contextual}{Contextual approximation and the separate full-abstraction question}
For \(H\vdash M,N:t\), write \(M\sqsubseteq_o N\) when every
well-typed program context \(C[-]\) closing both terms at ground type
satisfies: \(C[M]\Downarrow V\) implies \(C[N]\Downarrow V\).
As in the source, context binders may close free variables of their
inserted terms; that capture is specified by the context, not confused
with capture-avoiding substitution. Contexts retain their binder names.
Operational equivalence \(M\simeq_o N\) is approximation both ways.
Denotational order means pointwise order in every \(H\)-environment.
The source calls equality-implies-operational-equivalence adequacy;
we derive that conclusion from the ground theorem. Full abstraction
would additionally require operational equivalence to imply denotational
equality \cite[Section~4.3, p.~34]{Gunter1994TypesAuthorOA}.
\end{hbdefinition}

\begin{hbtheorem}{thm:p10-oa-contextual-soundness}{proved=annex:p10.oa.contextual-soundness}{The contextual consequence actually supported by adequacy}
Pointwise denotational order implies contextual approximation.
In particular denotational equality implies operational equivalence
for open as well as closed same-typed terms. Thus the constructed
interpretation is adequate in the source's contextual sense.
No converse is inferred from compositionality or from this implication.
\end{hbtheorem}

\begin{hbtheorem}{thm:p10-oa-strategy-boundaries}{proved=annex:p10.oa.strategy-boundaries}{Call-by-value and higher-type termination are different claims}
Let \(\Omega_t=\mu x:t.x\).
In this call-by-name language \((\lambda x:\mathsf{num}.0)\Omega_{\mathsf{num}}
\Downarrow0\), although \(\Omega_{\mathsf{num}}\Uparrow\).
Changing only application to call-by-value, by requiring its argument
to evaluate to a value before substitution, makes that application
diverge; the unchanged denotation remains zero and is therefore not
adequate for that changed strategy.
Also \(\lambda x:\mathsf{num}.\Omega_{\mathsf{num}}\) is a value
whose denotation is the bottom function, whereas
\(\Omega_{\mathsf{num}\to\mathsf{num}}\) diverges with that same
denotation. Higher-type termination is not a ground observation and
does not satisfy the displayed ground bottom criterion.
\end{hbtheorem}

\begin{hbtheorem}{thm:p10-oa-countdown}{proved=annex:p10.oa.countdown}{A recursive PCF computation, with syntax and approximants}
As a worked abbreviation, not a new primitive, put
\[
 K=\mu f:\mathsf{num}\to\mathsf{num}.
    \lambda x:\mathsf{num}.
    \mathsf{if}\ \operatorname{zero?}(x)\ \mathsf{then}\ 0\
    \mathsf{else}\ \operatorname{succ}(f(\operatorname{pred}(x))).
\]
For every closed \(P:\mathsf{num}\) with \(P\Downarrow\bar n\),
\(K P\Downarrow\bar n\); in particular \(K\bar3\Downarrow\bar3\).
Its semantic functional has iterates from bottom whose \(i\)-th
approximation gives \(n\) on precisely \(n<i\), and bottom on the
remaining inputs. The least fixed point is the strict identity on
\(\mathbb N_{\bot}\). Therefore \(K\Omega_{\mathsf{num}}\Uparrow\),
unlike the constant-zero application of the preceding theorem.
The operational proof permits unevaluated argument expressions \(P\),
not just already evaluated numeral arguments.
\end{hbtheorem}

\subsection{An adequate model that is not fully abstract}

\begin{hbdefinition}{def:p10-oa-sequential-relations}{Semantic relations for the sequentiality dependency}
At ground type let
\(Q_g=\{(a,b,c):a=\bot\text{ or }b=\bot\text{ or }a=b=c\}\).
At arrows a function triple is in \(Q_{s\to t}\) when it sends every
\(Q_s\)-related argument triple to a \(Q_t\)-related result triple.
Put \(A_t=\{d:(d,d,d)\in Q_t\}\). These selected elements are not
claimed to form a new full domain category. On \(A_g\), let \(B_g\)
be equality; put \((f,h)\in B_{s\to t}\) when both endpoints are
in \(A_{s\to t}\) and \((f(a),h(b))\in B_t\) whenever
\((a,b)\in B_s\). These editorial semantic logical relations prove
sequentiality; neither replaces the adequacy relation \(\mathcal R_t\).
\end{hbdefinition}

\begin{hbtheorem}{thm:p10-oa-sequential-relations}{proved=annex:p10.oa.sequential-relations}{Two complete semantic fundamental lemmas}
\(Q_t\) and \(B_t\) contain their constantly-bottom tuples and
are closed under directed suprema in the product carriers. Any triple
whose first or second coordinate is bottom belongs to \(Q_t\).
PCF interpretations preserve both relations on related environments.
Every closed denotation therefore belongs to \(A_t\), and inserting
\(B_t\)-related closed denotations into any typed closing ground context
produces equal ground denotations. This is a theorem about the source
syntax, not a premise that every continuous function is sequential.
\end{hbtheorem}

\begin{hbtheorem}{thm:p10-oa-not-fully-abstract}{proved=annex:p10.oa.not-fully-abstract}{The source's T,F counterexample with its sequentiality proved}
Put \(\Omega=\Omega_{\mathsf{bool}}\). Define
\[
 \begin{split}
 T={}&\lambda f:\mathsf{bool}\to\mathsf{bool}\to\mathsf{bool}.
 \mathsf{if}\ f\,\mathsf{true}\,\Omega\ \mathsf{then}\\
 &\quad\mathsf{if}\ f\,\Omega\,\mathsf{true}\ \mathsf{then}\\
 &\qquad\mathsf{if}\ f\,\mathsf{false}\,\mathsf{false}
    \ \mathsf{then}\ \Omega\ \mathsf{else}\ \mathsf{true}\\
 &\quad\mathsf{else}\ \Omega\ \mathsf{else}\ \Omega.
 \end{split}
\]
Let \(F\) be the identical term with the innermost else-result false
instead of true. Both have type
\((\mathsf{bool}\to\mathsf{bool}\to\mathsf{bool})\to\mathsf{bool}\).
Then \(T\simeq_o F\), although their denotations are different.
The separating argument is the curried continuous parallel disjunction
with the source truth table
\[
 \begin{array}{c|ccc}
 \operatorname{por}&\mathsf{true}&\mathsf{false}&\bot\\\hline
 \mathsf{true}&\mathsf{true}&\mathsf{true}&\mathsf{true}\\
 \mathsf{false}&\mathsf{true}&\mathsf{false}&\bot\\
 \bot&\mathsf{true}&\bot&\bot
 \end{array}
\]
It gives \(\llbracket T\rrbracket(\operatorname{por})=\mathsf{true}\)
and \(\llbracket F\rrbracket(\operatorname{por})=\mathsf{false}\).
Thus the continuous model is not fully abstract for sequential PCF.
The contextual equivalence is proved directly here without importing
stable-model adequacy \cite[Tables~9--10, pp.~36--38]{Gunter1994TypesAuthorOA}
\cite[pp.~234--236]{Plotkin1977LCFAuthorOA}.
\end{hbtheorem}

\begin{hbmethod}{def:p10-oa-use}{Applying the adequacy theorem}
Fix the source language and evaluation strategy; check typing and identify
the ground numeral/boolean observation. Interpret recursion by the least
continuous fixed point, prove substitution and evaluation soundness,
and verify bottom/closure of the typed semantic--syntactic relation.
Use its fundamental lemma for the converse observation direction.
For contextual claims quantify over the actual typed closing contexts.
Changing evaluation to call-by-value, adding effects or parallel operations,
or treating function termination as an observation changes the proof target.
\end{hbmethod}

\begin{proofboundary}{def:p10-oa-source-boundary}{Source-specific scope and stronger language variants}
This supplies the omitted proof for Gunter's continuous bc-domain
call-by-name PCF interpretation, not its stable dI-domain interpretation
or all later record/subtyping languages. The predecessor-zero behavior
differs from Plotkin's original arithmetic constants. The continuous
model's source full-abstraction counterexample now has its sequentiality
dependency proved above. The stronger parallel-extension full-abstraction
and definability theorems do not follow from that failure proof.
Abramsky--Jung's predomain/partial-map and monadic
call-by-value proposals need their own value/computation semantics and
adequacy argument. Printed chapters, the held author presentations and
the whole-source allocations remain distinct.
\end{proofboundary}
